\documentclass[a4paper,10pt]{article}
%% Generated by _tools/cv.py from _src/site.toml. Edit the data there, then rerun.
\usepackage[utf8]{inputenc}
\usepackage[T1]{fontenc}
\usepackage{lmodern}
\usepackage{fontawesome5}
\usepackage{microtype}
\usepackage{textcomp}
\usepackage{geometry}
\usepackage{enumitem}
\usepackage{titlesec}
\usepackage{xcolor}
\usepackage{hyperref}

% Real Unicode text layer, so applicant-tracking systems and copy-paste read
% ligatures and accents correctly.
\input{glyphtounicode}
\pdfgentounicode=1

% Black text; navy for headings and links only.
\definecolor{ink}{HTML}{000000}
\definecolor{accent}{RGB}{50,50,100}
\definecolor{muted}{HTML}{555555}

\geometry{top=0.34in, bottom=0.32in, left=0.6in, right=0.6in}

\color{ink}
\pagestyle{empty}
\setlength{\parindent}{0pt}
\setlength{\parskip}{1pt}

\titleformat{\section}{\normalfont\large\bfseries\color{accent}}{}{0em}{}[{\color{muted}\titlerule[0.4pt]}]
\titlespacing*{\section}{0pt}{5pt}{3pt}

\setlist[itemize]{noitemsep, topsep=1pt, leftmargin=1.2em, label=\textbullet}

\hypersetup{
  colorlinks=true, linkcolor=accent, urlcolor=accent,
  pdftitle={Mike Silva -- Curriculum Vitae},
  pdfauthor={Mike Silva},
  pdfkeywords={computational physics, complex systems, nuclear reactions, reduced-order models, stochastic processes}
}

% \entry{Institution}{City, Country}{Title / degree}{Dates}
\newcommand{\entry}[4]{%
  \textbf{#1}\hfill #4\\
  \textit{#3}\hfill {\color{muted}#2}\par}
\newcommand{\coursework}[1]{{\small\color{muted}Coursework: #1}\par\vspace{2pt}}
\newcommand{\skill}[2]{\hangindent=9.5em\hangafter=1\makebox[9.5em][l]{\textbf{#1}}#2\par}

\begin{document}

% ---- header: no phone, no photo, no birth date (international norm; the PDF is public)
\begin{center}
  {\LARGE\textbf{Mike Silva}}\\[4pt]
  \small
  \newcommand{\ico}[1]{{\color{muted}#1}~}%
  \ico{\faMapMarker*}Paris, France \qquad
  \href{mailto:pro.mikesilva@gmail.com}{\ico{\faEnvelope}pro.mikesilva@gmail.com} \qquad
  \href{https://mksv-pro.github.io}{\ico{\faGlobe}mksv-pro.github.io} \qquad
  \href{https://github.com/mksv-pro}{\ico{\faGithub}mksv-pro} \qquad
  \href{https://www.linkedin.com/in/mkesilva/}{\ico{\faLinkedin}mkesilva}
\end{center}

\textbf{Research interests:} complex systems~\textperiodcentered{} quantum information~\textperiodcentered{} stochastic processes~\textperiodcentered{} machine learning for physics~\textperiodcentered{} agent-based and mean-field models~\textperiodcentered{} quantitative finance.

% ------------------------------------------------------------------------------
\section*{Education}

\entry{Université Paris-Saclay}{Orsay, France}{MSc Physics of Complex Systems}{2026 -- 2027 (expected)}
\coursework{stochastic processes, statistical field theory, disordered systems, random matrices, machine learning.}

\entry{Sorbonne Université}{Paris, France}{MSc Quantum Information}{2026 -- 2027 (expected)}
\coursework{quantum algorithms, quantum information theory, quantum optics, quantum metrology, ultracold atoms.}

\entry{Université Paris Dauphine-PSL}{Paris, France}{MSc Economics \& Finance; BSc Economics \& Finance, third year (L3, 2025–2026)}{2025 -- 2027 (expected)}
\coursework{derivatives, portfolio management, time-series econometrics, macroeconomics.}

\entry{Université Paris Cité}{Paris, France}{MSc Fundamental Physics, first year, \textit{mention Bien}}{2025 -- 2026}
\coursework{statistical, high-energy, computational and nonlinear physics, quantum information, condensed matter physics.}

\entry{Sorbonne Université}{Paris, France}{MSc Mathematics, first year; BSc Physics and Mathematics (double major)}{2021 -- 2025}
\coursework{probability, stochastic control, jump processes, functional analysis, PDEs, numerical methods, statistical physics.}

\entry{Université Paris 1 Panthéon-Sorbonne}{Paris, France}{BA Economics and Political Science}{2022 -- 2025}
\coursework{econometric methods, macro- and microeconomics, monetary economics, political economy.}

% ------------------------------------------------------------------------------
\section*{Research Experience}

\entry{IJCLab – Irène Joliot-Curie Laboratory (CNRS / Université Paris Cité)}{Orsay, France}%
{Research Intern, supervised by Chloé Hebborn and Patrick McGlynn}{Jun 2026 -- Jul 2026}
Emulation of nuclear reaction models for uncertainty quantification on optical potentials.
Thesis: \textit{Computational Emulation of Nuclear Dynamics}, graded 18/20
(\href{https://github.com/mksv-pro/Mike-M1-Internship-IJCLab-2026}{report and code}).
\begin{itemize}
  \item Developed a coupled-channel scattering solver based on the Direct Boundary Matching Method (DBMM), validated against an independent R-matrix calculation
  \item Built two emulators (reduced basis method; learning reduced-order model, first coupled-channel extension) matching the solver within 5\,\texttimes\,10\textsuperscript{\textminus{}3} (median) on four systems with up to 8 coupled channels and 10 parameters
  \item Benchmarked both against the solver: speed-ups grow with system size, and the learning model gains 1.3–1.9\texttimes{} more
\end{itemize}
\vspace{3pt}

\entry{CEREMADE, Université Paris Dauphine-PSL}{Paris, France}%
{Visiting Research Student, supervised by Pierre Cardaliaguet}{Sep 2024 -- Jan 2025}
Mean-field games applied to the propagation of monetary shocks in a general-equilibrium economy with sticky prices.
\begin{itemize}
  \item Studied differential games and mean-field models of strategic interactions between heterogeneous agents
  \item Analyzed the coupled Hamilton–Jacobi–Bellman and Fokker–Planck equations arising in dynamic economic models
  \item Applied mathematical and computational tools to monetary policy and collective decision-making
\end{itemize}
\vspace{3pt}

\entry{Sciences Po}{Paris, France}%
{Research Assistant, supervised by Julia Cagé}{Jun 2023 -- Aug 2023}
Research assistant for \textit{Hosting Media Bias} (Cagé et al.): built part of its dataset of French broadcasts.

% ------------------------------------------------------------------------------
\section*{Research Projects}

\entry{Urban morphogenesis: diffusion-limited aggregation and Schelling's model}{Paris, France}%
{Université Paris Cité, supervised by Raphael Raynaud}{Jan 2026}
Agent-based model coupling Schelling segregation with diffusion-limited aggregation: emergent morphologies and fractal dimension; extension with reinforcement learning under way.
\vspace{3pt}

\entry{Self-gravitating N-body systems}{Paris, France}%
{Sorbonne Université, supervised by Lionel Foret}{Feb 2024}
Direct-integration simulation of gravitational N-body dynamics and long-range collective effects.

% ------------------------------------------------------------------------------
\section*{Skills}

\skill{Programming}{Python (NumPy, SciPy, Numba, JAX, PyTorch), R; Git, pytest, CI, LaTeX}
\skill{Numerics \& ML}{PDE solvers, reduced-order models, Monte Carlo, neural networks, reinforcement learning}
\skill{Quantum}{quantum algorithms, quantum information theory, spin glasses, random matrices}
\skill{Complex systems}{agent-based models, mean-field games, stochastic processes, fractal growth}
\skill{Economics}{stochastic calculus, derivatives pricing, econometrics, general equilibrium}
\skill{Languages}{French (native), English (fluent), Portuguese (basic), German (basic)}

\end{document}
